\documentclass{article}
\usepackage{fontspec}

\setmainfont{Libertinus Serif}

\begin{document}

Приведение $\lambda$-терма

Рассмотрим терм

\[((\lambda a.(\lambda b.b\,b)(\lambda b.b\,b))b)\
((\lambda c.(c\, b)) (\lambda a.a))\]

Выполним последовательные $\beta$-редукции.

Шаг 1. Редуцируем левую часть:

\[((\lambda a.(\lambda b.b\,b)(\lambda b.b\,b))b)\to
(\lambda b.b\,b)(\lambda b.b\,b)\]

После этого весь терм принимает вид

\[(\lambda b.b\,b)(\lambda b.b\,b)
((\lambda c.(c\, b))(\lambda a.a))\]

Шаг 2. Редуцируем правую часть:

\[(\lambda c.(c\, b))(\lambda a.a)\to
(\lambda a.a)b\]

Получаем

\[(\lambda b.b\,b)(\lambda b.b\,b)
((\lambda a.a)\, b)\]

Шаг 3. Выполняем ещё одну $\beta$-редукцию:

\[(\lambda a.a)\, b \to b\]

Терм:

\[(\lambda b.b\,b)(\lambda b.b\,b) b\]

Шаг 4. Рассмотрим выражение
\[(\lambda b.b\,b)(\lambda b.b\,b)\]

После $\beta$-редукции получаем

\[(\lambda b.b\,b)(\lambda b.b\,b)\to
(\lambda b.b\,b)(\lambda b.b\,b)\]

то есть то же самое выражение.

Таким образом, подтерм
\[(\lambda b.b\,b)(\lambda b.b\,b)\]
редуцируется сам в себя и порождает бесконечную последовательность редукций.

Заключение:

Данный $\lambda$-терм не имеет нормальной формы, так как содержит
зацикливающийся подтерм (омега-комбинатор).

\end{document}